% Standalone LaTeX section for the revised tDWfPINNs manuscript.
% This file can be compiled as a standalone document, or the section body can be
% copied into the numerical-experiments section of the paper.

\documentclass[11pt]{article}
\usepackage[a4paper,margin=1in]{geometry}
\usepackage{amsmath,amssymb,amsfonts,bm}
\usepackage{graphicx}
\usepackage{booktabs}
\usepackage{hyperref}

\newcommand{\Dtalpha}{\partial_t^\alpha}
\newcommand{\R}{\mathbb{R}}
\newcommand{\bb}{\bm b}
\newcommand{\grad}{\nabla}

\begin{document}

\section{Two-dimensional nonlinear heterogeneous diffusion-wave equation on an irregular domain}
\label{sec:irregular-2d-dw}

To further assess the proposed transformed diffusion-wave fPINN formulation in a genuinely multidimensional and geometrically nontrivial setting, we consider a two-dimensional nonlinear heterogeneous time-fractional diffusion-wave equation on an irregular domain. This example is designed to test four aspects that are not visible in one-dimensional benchmark problems: multidimensionality, a non-rectangular geometry, spatially heterogeneous diffusion, and nonlinear interaction.

Let $T=1$ and let $1<\alpha<2$. The spatial domain is a square with an off-center circular hole,
\begin{equation}
\Omega=
(-1,1)^2\setminus \overline{B_{r_0}(c)},
\qquad
c=(-0.3,0.2),
\qquad
r_0=0.25,
\label{eq:irregular-domain}
\end{equation}
where
\[
B_{r_0}(c)=\left\{(x,y)\in\R^2:(x+0.3)^2+(y-0.2)^2<r_0^2\right\}.
\]
We solve
\begin{equation}
\Dtalpha u
-
\grad\cdot\left(a(x,y)\grad u\right)
+
\bb(x,y)\cdot\grad u
+
\lambda u^3
=
f(t,x,y),
\qquad
(t,x,y)\in(0,T]\times\Omega,
\label{eq:irregular-2d-pde}
\end{equation}
with homogeneous Dirichlet boundary conditions and homogeneous initial conditions,
\begin{equation}
 u(t,x,y)=0,
\qquad
(t,x,y)\in(0,T]\times\partial\Omega,
\label{eq:irregular-2d-bc}
\end{equation}
\begin{equation}
 u(0,x,y)=0,
\qquad
 \partial_t u(0,x,y)=0,
\qquad
(x,y)\in\Omega.
\label{eq:irregular-2d-ic}
\end{equation}
The heterogeneous diffusion coefficient, convection field, and nonlinear coefficient are chosen as
\begin{equation}
 a(x,y)=1+0.3\sin(\pi x)\cos(\pi y),
\qquad
 \bb(x,y)=\begin{pmatrix}1+y\\ x-1\end{pmatrix},
\qquad
 \lambda=1.
\label{eq:irregular-2d-coefficients}
\end{equation}
The diffusion coefficient satisfies
\[
0.7\leq a(x,y)\leq 1.3,
\]
so the elliptic part remains uniformly positive.

\subsection{Manufactured exact solution and source term}

We use the manufactured exact solution
\begin{equation}
 u^\ast(t,x,y)=q(t)\phi(x,y),
\qquad
 q(t)=t^2(1-t)^2,
\label{eq:irregular-2d-exact}
\end{equation}
where
\begin{equation}
 \phi(x,y)
 =
 (1-x^2)(1-y^2)
 \left((x+0.3)^2+(y-0.2)^2-r_0^2\right).
\label{eq:irregular-2d-phi}
\end{equation}
This construction enforces homogeneous boundary conditions on both parts of the boundary. Indeed, $(1-x^2)(1-y^2)=0$ on the outer square boundary $x=\pm1$ or $y=\pm1$, while
\[
(x+0.3)^2+(y-0.2)^2-r_0^2=0
\]
on the inner circular boundary. Moreover, since $q(0)=q'(0)=0$, the initial conditions in \eqref{eq:irregular-2d-ic} are also satisfied.

The temporal Caputo derivative is available in closed form. Since
\[
q(t)=t^2-2t^3+t^4,
\]
we have, for $1<\alpha<2$,
\begin{equation}
 \Dtalpha q(t)
 =
 \frac{2}{\Gamma(3-\alpha)}t^{2-\alpha}
 -
 \frac{12}{\Gamma(4-\alpha)}t^{3-\alpha}
 +
 \frac{24}{\Gamma(5-\alpha)}t^{4-\alpha}.
\label{eq:irregular-2d-caputo-q}
\end{equation}
Thus
\begin{equation}
 \Dtalpha u^\ast(t,x,y)=\left(\Dtalpha q(t)\right)\phi(x,y).
\label{eq:irregular-2d-caputo-u}
\end{equation}
The source term is defined by substituting the manufactured solution into \eqref{eq:irregular-2d-pde}:
\begin{equation}
 f(t,x,y)
 =
 \Dtalpha u^\ast
 -
 \grad\cdot\left(a\grad u^\ast\right)
 +
 \bb\cdot\grad u^\ast
 +
 \lambda (u^\ast)^3.
\label{eq:irregular-2d-source-def}
\end{equation}
Equivalently, using the separable form $u^\ast=q\phi$, this can be written compactly as
\begin{equation}
\begin{aligned}
 f(t,x,y)
 ={}&
 \left(\Dtalpha q(t)\right)\phi(x,y)
 -
 q(t)\left[\grad a(x,y)\cdot\grad\phi(x,y)+a(x,y)\Delta\phi(x,y)\right]
 \\
&
 +
 q(t)\,\bb(x,y)\cdot\grad\phi(x,y)
 +
 \lambda q(t)^3\phi(x,y)^3.
\end{aligned}
\label{eq:irregular-2d-source-compact}
\end{equation}
For implementation, it is not necessary to expand the polynomial $\phi$. A stable and transparent way to evaluate the source is to write
\begin{equation}
 \phi(x,y)=\psi(x,y)\chi(x,y),
 \qquad
 \psi(x,y)=(1-x^2)(1-y^2),
 \qquad
 \chi(x,y)=(x+0.3)^2+(y-0.2)^2-r_0^2.
\label{eq:irregular-2d-psi-chi}
\end{equation}
Then
\begin{equation}
 \grad\phi=\chi\grad\psi+\psi\grad\chi,
 \qquad
 \Delta\phi=\chi\Delta\psi+2\grad\psi\cdot\grad\chi+\psi\Delta\chi,
\label{eq:irregular-2d-phi-derivatives}
\end{equation}
where
\begin{equation}
\begin{aligned}
 \grad\psi
 &=
 \begin{pmatrix}
 -2x(1-y^2)\\
 -2y(1-x^2)
 \end{pmatrix},
&
 \Delta\psi
 &=2x^2+2y^2-4,\\
 \grad\chi
 &=
 \begin{pmatrix}
 2(x+0.3)\\
 2(y-0.2)
 \end{pmatrix},
&
 \Delta\chi
 &=4,
\end{aligned}
\label{eq:irregular-2d-aux-derivatives}
\end{equation}
and
\begin{equation}
 \grad a(x,y)
 =
 \begin{pmatrix}
 0.3\pi\cos(\pi x)\cos(\pi y)\\
 -0.3\pi\sin(\pi x)\sin(\pi y)
 \end{pmatrix}.
\label{eq:irregular-2d-grad-a}
\end{equation}
Equations \eqref{eq:irregular-2d-source-compact}--\eqref{eq:irregular-2d-grad-a} give a complete analytic expression for the manufactured forcing term.

\subsection{PINN setup for the irregular-domain benchmark}

The loss function follows the same structure as in the preceding experiments. The interior residual is evaluated at collocation points sampled from $(0,T]\times\Omega$, where spatial points are generated by rejection sampling in $(-1,1)^2$ and discarding points inside $B_{r_0}(c)$. Boundary points are sampled on both the outer square boundary and the inner circular boundary. For the circular boundary, we use the parametrization
\begin{equation}
 x=-0.3+r_0\cos\theta,
 \qquad
 y=0.2+r_0\sin\theta,
 \qquad
 \theta\in[0,2\pi).
\label{eq:irregular-2d-circle-param}
\end{equation}
The initial points are sampled inside $\Omega$ at $t=0$.

In this benchmark, the primary objective is not to rescan all quadrature configurations, but to demonstrate that the proposed transformed formulation remains effective in a two-dimensional nonlinear problem with heterogeneous coefficients and an irregular geometry. We therefore report results for the transformed schemes, especially GJ-II and MC-II, and optionally include GJ-I as a direct-form reference. A typical setting is
\begin{equation}
 \alpha\in\{1.25,1.5,1.75\},
 \qquad
 M_{\rm GJ}\in\{32,64\},
 \qquad
 M_{\rm MC}\in\{640,1280\}.
\label{eq:irregular-2d-parameters}
\end{equation}
The relative error is evaluated on a separate test set sampled from $(0,T]\times\Omega$:
\begin{equation}
 e_r(\theta)
 =
 \frac{
 \left(\sum_{i=1}^{N_{\rm test}}|u_\theta(t_i,x_i,y_i)-u^\ast(t_i,x_i,y_i)|^2\right)^{1/2}
 }{
 \left(\sum_{i=1}^{N_{\rm test}}|u^\ast(t_i,x_i,y_i)|^2\right)^{1/2}
 }.
\label{eq:irregular-2d-relative-error}
\end{equation}

\subsection{Suggested figures and table}

For this experiment, we recommend reporting: (i) the irregular domain and sampled collocation points; (ii) the exact solution, the neural-network prediction, and the absolute error at a representative time such as $t=0.5$; and (iii) a table summarizing the relative error, training time, and GPU memory consumption for different $\alpha$ and quadrature strategies.

\begin{table}[t]
\centering
\caption{Suggested result table for the two-dimensional nonlinear heterogeneous diffusion-wave equation on the irregular domain. The entries should be filled after training.}
\label{tab:irregular-2d-results}
\begin{tabular}{ccccc}
\toprule
$\alpha$ & Method & $M$ & Relative $L^2$ error & Time / memory \\
\midrule
1.25 & MC-II & 640 or 1280 & -- & -- \\
1.25 & GJ-II & 32 or 64 & -- & -- \\
1.50 & MC-II & 640 or 1280 & -- & -- \\
1.50 & GJ-II & 32 or 64 & -- & -- \\
1.75 & MC-II & 640 or 1280 & -- & -- \\
1.75 & GJ-II & 32 or 64 & -- & -- \\
\bottomrule
\end{tabular}
\end{table}

\paragraph{Remark on the domain assumption.}
This example is posed on a bounded Lipschitz domain with a smooth inner boundary. If the general problem statement in the manuscript currently restricts $\Omega$ to convex polyhedral domains, it should be relaxed to bounded Lipschitz domains, or this example should be explicitly described as an irregular-domain extension of the main computational framework.

\end{document}
